把一张图片的每个像素改动一个肉眼看不出的量，一个准确率 \(96\%\) 的分类器会掉到个位数。同样幅度的随机噪声几乎不影响它。

这两句话放在一起，就把这个单元要讲的事情框住了。模型不是对噪声敏感——它对噪声很不敏感。真正的情况是：在高维输入空间里存在少数特殊方向，沿着它们走一小步就能越过决策边界，而这些方向不难找到。

根子在几何。一个 \(\ell_\infty\) 球的顶点到中心的欧氏距离是它半径的 \(\sqrt d\) 倍，\(d\) 是几千的时候这个倍数是几十；人对``改动很小''的判断按 \(\ell_\infty\) 建立，模型的决策依据却是内积。同一份数据，两把尺子量出来的结果差了一个数量级不止。与此同时，训练只在有限个点上验证过模型的行为，而这些点的 \(\varepsilon\)-邻域并集在整个空间中的体积占比小到可以忽略不计。

于是鲁棒性被写成了一个关于邻域的全称命题——对每个输入、对邻域内每个扰动，预测都不变。这个形状与上一个单元的隐私定义如出一辙，也带来同一个后果：它不能靠采样来验证。跑一万次攻击都没打穿，不构成任何保证。

\hyperref[sec:17-Constrained-Guarantees/02-Adversarial-Robustness/01]{``微小扰动为何能翻转预测''}把线性化的解释算清楚：\(\ell_\infty\) 约束下线性泛函的最大变化由 \(\ell_1\) 范数给出，随维数线性增长，而随机方向只能拿到 \(\sqrt d\) 分之一。同时把高维几何的两条计算摆出来，说明``在数据点上正确''与``在数据点邻域上正确''相差多远。

\hyperref[sec:17-Constrained-Guarantees/02-Adversarial-Robustness/02]{``Lipschitz 常数就是认证半径''}给出唯一算数的那类论证。间隔函数的 Lipschitz 常数直接控制一个不依赖攻击方法的半径，证明只有一行；难处在于这个常数精确计算是 NP 难的，可算的上界又极其松。随机平滑换了条路，用高斯卷积把 Lipschitz 性质制造出来，换来的是推理时要采样上千次、结论带一个显著性水平。

\hyperref[sec:17-Constrained-Guarantees/02-Adversarial-Robustness/03]{``鲁棒与准确不可兼得''}给出这个单元的硬结论：存在数据分布，使任何分类器的标准准确率与鲁棒准确率之和被一个小于 2 的常数卡住。对抗训练是在这条约束线上选点，而它优化的是真实目标的一个下界的近似——这一点决定了为什么梯度遮蔽类的``防御''能在评测里好看，又在换一种攻击时全线失效。

读完带走一条判断就够：攻击失败不构成安全性证明。可证明的界哪怕数字难看也是承诺，攻击失败的记录哪怕数字漂亮也只是观察，两者不能互相替代，也不能互相佐证。
